% Generated draft: requires mathematical and source review.
\documentclass[11pt]{article}
\usepackage[margin=25mm]{geometry}
\usepackage{fontspec}
\setmainfont{Latin Modern Roman}
\usepackage{amsmath,amssymb,mathtools,unicode-math}
\setmathfont{Latin Modern Math}
\usepackage{enumitem,needspace,xurl,hyperref}
\setlength{\parindent}{0pt}
\setlength{\parskip}{5pt}
\setlength{\emergencystretch}{3em}
\newcommand{\N}{\mathbb{N}}
\newcommand{\Z}{\mathbb{Z}}
\newcommand{\Q}{\mathbb{Q}}
\newcommand{\R}{\mathbb{R}}
\newcommand{\C}{\mathbb{C}}
\newcommand{\F}{\mathbb{F}}
\newcommand{\A}{\mathbb{A}}
\newcommand{\PP}{\mathbb{P}}
\newcommand{\E}{\mathbb{E}}
\newcommand{\eps}{\varepsilon}
\newcommand{\dd}{\,\mathrm{d}}
\begin{document}
\section*{UnsolvedMath ID: 30006166}
\textbf{Source problem code:} catalogue entry
\subsection{Definitions and mathematical statement}
Let \(X=2^{\N}\) be Cantor space with its product topology. With the usual restricted-wreath-product convention, put
\[
G=\Z\wr\Z=\left(\bigoplus_{k\in\Z}\Z\right)\rtimes\Z.
\]
Here \(\bigoplus_{k\in\Z}\Z\) consists of the finitely supported functions \(f:\Z\to\Z\), and \(m\in\Z\) acts by \((\sigma_m f)(k)=f(k-m)\). Thus
\[
(f,m)(h,n)=(f+\sigma_m h,m+n).
\]
A continuous action of \(G\) on \(X\) is a homomorphism \(\alpha:G\to\operatorname{Homeo}(X)\). Its orbit equivalence relation is
\[
x\mathrel{E_\alpha}y\quad\Longleftrightarrow\quad (\exists g\in G)\;\alpha(g)(x)=y.
\]
The relation \(E_\alpha\) is \emph{hyperfinite} if there are Borel equivalence relations \(F_n\subseteq X\times X\), indexed by \(n\in\N\), such that every \(F_n\)-class is finite and
\[
F_n\subseteq F_{n+1}\quad\text{for every }n\in\N,
\qquad E_\alpha=\bigcup_{n\in\N}F_n.
\]
This is a Borel condition on the entire space \(X\); no measure or exceptional null set is specified.

\textbf{Source question.} Is the generic action of \(G=\Z\wr\Z\) on Cantor space hyperfinite?

\textbf{Unresolved meaning of genericity.} The supplied extract does not define ``the generic action'' or its ambient action space. A conventional Baire-category interpretation uses
\[
\mathcal A=\operatorname{Hom}\bigl(G,\operatorname{Homeo}(X)\bigr),
\]
with the subspace topology inherited from \(\operatorname{Homeo}(X)^G\), where \(\operatorname{Homeo}(X)\) has the compact-open topology. Under this interpretation the precise question is whether there exists a comeager set \(C\subseteq\mathcal A\) such that
\[
\forall\alpha\in C\;\exists(F_n)_{n\in\N}\quad
\left[
\begin{gathered}
\text{each }F_n\text{ is a Borel equivalence relation on }X
\text{ with finite classes},\\
F_n\subseteq F_{n+1}\text{ for every }n\in\N,
\qquad E_\alpha=\bigcup_{n\in\N}F_n
\end{gathered}
\right].
\]
A set is comeager when its complement is a countable union of nowhere dense sets. This displayed formulation is conditional on the stated interpretation. If the source instead singles out an action whose conjugacy class is comeager, the question concerns that action's orbit relation; existence of such a conjugacy class is not established by the supplied extract. The extract also does not specify whether genericity is relative to all actions or to a restricted space such as the free actions. These qualifications require source clarification before the conditional formulation can be treated as definitive.

The surrounding source discussion places this example among unresolved questions about generic actions, including actions of amenable groups, after a result called Theorem~1. It attributes that theorem's proof to a connection between finite Borel asymptotic dimension and hyperfiniteness developed in its reference~[1]. Neither Theorem~1 nor the bibliographic identification of that numbered reference is included in the extract; no additional hypothesis or conclusion from them is imposed here.
\subsection{Short English statement}
The question asks whether the orbit equivalence relation of a generic continuous action of \(\Z\wr\Z\) on Cantor space can be written as an increasing union of Borel equivalence relations with finite classes. The supplied extract leaves the meaning and ambient space of ``generic'' unspecified, so the detailed Baire-category formulation above remains conditional. The supplied record labels the problem open; this draft does not independently verify that status.
\subsection{Sources}
\begin{itemize}
\item Local UnsolvedMath registry: supplied complete source record for permanent ID \(30006166\), entitled \emph{Hyperfiniteness of Generic Wreath-Product Actions}; exact source problem code \texttt{catalogue entry}. The record assigns the problem to the Oberwolfach Reports Open Problems collection and gives proposed year \(2025\); no individual proposer is supplied.
\item Supplied primary citation: \emph{Set Theory} (2025), Oberwolfach Reports, pp.~75--140, DOI \texttt{10.4171/owr/2025/2}. \url{https://doi.org/10.4171/owr/2025/2}. The supplied provenance describes Oberwolfach Reports as records of mathematical workshops and their open-problem discussions.
\item Supplied literature source: \emph{Oberwolfach Report 2/2025}. \url{https://content.ems.press/assets/public/full-texts/serials/owr/22/1/14299082/online/10.4171-owr-2025-2.pdf}.
\item Supplied literature source: Iyer--Shinko, \emph{Asymptotic dimension and hyperfiniteness}. \url{https://arxiv.org/abs/2409.03078}.
\item Supplied dated literature curation, checked \texttt{2026-08-22}: the record labels the problem \emph{open} and states that the primary report explicitly lists this generic-action question as open and that known finite-asymptotic-dimension results do not cover this group. This is supplied dated curation, not an independently verified literature assessment. The background expressly requests independent verification of time-sensitive status and cited claims.
\item Supplied statement-verification metadata: \texttt{corrected\_verified}; the record says that the statement was reconstructed and checked against the official Oberwolfach report. This reports the registry's verification claim and does not constitute independent verification by this draft.
\end{itemize}
\end{document}
